\documentclass[11pt,a4paper]{article}
\usepackage{kotex}
\usepackage{amsmath,amssymb,booktabs,geometry}
\geometry{margin=25mm}
\title{I01 선검증 1: 위상 표지와 유한 시계 아래 관측 rank의 경계}
\author{RELATIONAL ATLAS}
\date{2026년 10월 4일}
\begin{document}
\maketitle

\section{판정}
이 실험은 새로운 일반 이론의 확증이 아니다. 실험 전에 고정한
\texttt{prevalidation\_protocol\_2026-10-04.tex}의 격자 안에서는
충분한 위상 표지($L=2n+1$)를 보존하면 비제약 상태의 관측 Fisher rank가 $n$,
표지를 지우면 $n-1$이라는 개입 패턴이 유지됐다. 그러나 짧은 격자
($L=2n-1$)에서는 표지를 지워도 공통 위상 정보가 남는 반례가 나왔다.
따라서 ``표지 삭제는 언제나 한 차원을 지운다''는 명제는 \emph{기각}하고,
``충분한 표지 해상도와 이 검출기족에서만''이라는 조건부 명제로 남긴다.
이는 공간 차원의 동역학적 선택이나 세 방향의 특권성을 뜻하지 않는다.

\section{사전 고정한 설계}
$k_i\in\{-1,0,1\}$, $E(k)=\sum_i k_i^2$,
$w_k\propto e^{-\beta E(k)}$에서 $P=\sum_i k_i=0$인 상태와 비제약 상태를
따로 준비했다. 축 $i$의 복귀 검출기 바이어스 $b$, 기록된 공통 위상
$\phi_t=2\pi t/L$, 시계 위상 $\tau$를 사용했다. 사전 고정 격자는
\begin{equation}
n=3,4,5,6;\quad \beta=0.45,1.20;\quad b=0.55,1.25;\quad
\tau/\pi=0,1/2;\quad L=2n-1,2n+1.
\end{equation}
128개의 상태--조건 행을 계산하고, 각 행에서 같은 표본의 표지 보존/삭제를
비교했다. 조건부 복귀 확률은
\begin{equation}
q_i(t;\boldsymbol\theta,\tau)
=\left|\sum_k w_k\exp\!\left[i\left(E(k)\tau+
\sum_j k_j(\theta_j+\phi_t-b\delta_{ij})\right)\right]\right|^2
\end{equation}
이다. 표지 보존은 $q_i(t)$별 Bernoulli 정보를 합치고, 표지 삭제는
$\bar q_i=L^{-1}\sum_t q_i(t)$만 사용한다. 보존/삭제는 새로운 물리 제약의
삽입이 아니라 기록의 사용 여부를 바꾼 분석 개입이다.

\section{정확 확률과 회로 대조}
\begin{center}
\begin{tabular}{@{}lccc@{}}
\toprule
조건 & 상태 QFI rank & 표지 보존 CFI rank & 표지 삭제 CFI rank \\
\midrule
비제약, $L=2n+1$ & $n$ & $n$ & $n-1$ \\
$P=0$, $L=2n+1$ & $n-1$ & $n-1$ & $n-1$ \\
\bottomrule
\end{tabular}
\end{center}
사전 격자의 충분한 표지 64행에서 이 패턴을 벗어난 행은 0개였다.
상태 QFI의 $n/(n-1)$ 패턴은 전체 128행에서 실패 0개였다.
이 행들은 동일한 상태를 다른 측정 조건에서 반복 사용한 결정론적 격자이며,
64회 또는 128회의 독립 통계 시행으로 해석하지 않는다.
확률 미분의 독립 중앙차분 대조 최대 오차는 $3.38\times10^{-11}$이었다.
선택한 $n=3,4$ 회로 네 개의 Statevector와 수식의 결합확률 최대 차이는
$6.52\times10^{-14}$, $n=3$ 두 회로의 Aer density-matrix 대조 최대 차이는
$3.13\times10^{-15}$였다. 이것은 구현 검증이지 가족 전체의 독립 증명은 아니다.

짧은 표지 $L=2n-1$인 비제약 32행 중 16행에서, 표지 삭제 뒤에도 공통 위상
$0\to0.23$ rad 이동이 확률을 $10^{-10}$보다 크게 바꿨다. 16행 모두
$\tau=0$이고 $\tau=\pi/2$의 나머지 16행은 이 기준 아래 0이었다.
예컨대 $n=3,\beta=0.45,b=1.25,L=5,\tau=0$의 삭제 CFI rank는 3,
최대 확률 이동은 $0.00827$이다. 동일한 조건의 $\tau=\pi/2$에서는 삭제
rank가 2이고 공통 이동은 수치 정밀도에서 0이다.
짧은 격자에서는 허용된 전하 차이 $|P-P'|\le2n$ 중
$|P-P'|=2n-1$ 조화가 평균 후에도 살아남을 수 있다.
$\tau=\pi/2$에서의 소거는 이 에너지 스펙트럼의 간섭 특성으로 관찰됐으며,
다른 스펙트럼에 대한 정리가 아니다.
또한 $n=6,\beta=1.20,L=11,\tau=0$ 두 바이어스에서는
공통방향 정보가 약 $2\times10^{-12}$로 양수인데 상대 고유값
$10^{-8}$ 기준의 운영적 rank는 5다. 정확한 식별 가능성과 유한 감도
판정은 구분해야 한다.

\section{Aer 유한 샷에서의 오판 점검}
$n=3,\beta=0.45,b=1.25$, $L=5,7$, $\tau=0,\pi/2$의 비제약 및
$P=0$ 회로를 각 4096 shots로 실행했다. 각 상태 144회로, 합계 288회로다.
시계 후선택률은 비제약 0.25005, $P=0$ 0.24967로 설계값 $1/4$에 가까웠다.
공통 위상 블록 $0\to0.23$ rad의 2000회 재배분 검정 16개에서,
비제약/표지 보존 네 개의 $p$값은 $0.00050,0.00050,0.00150,0.00300$이었다.
비제약/표지 삭제 네 개는 $0.11294,0.78261,0.03548,0.03598$이었다.
마지막 둘은 정확 모형에서 공통방향 정보가 0인 $\tau=\pi/2$의
명목 5\% 거짓 양성이다. $P=0$ 여덟 검정 중 하나도 명목
$p=0.04748$이었으나 정확 모형은 공통 위상에 불변이다.
16개를 하나의 가족으로 본 Bonferroni 문턱 $0.003125$에서는
보존/비제약 네 개만 남고, 삭제 및 $P=0$은 모두 0개다.
\textbf{이 다중검정 보정은 raw $p$값을 본 뒤 추가한 탐색적 진단이며
사전 등록된 확증 검정으로 제시하지 않는다.}
반대로 $L=5,\tau=0$의 삭제 상태는 정확하게 rank 3이지만 이 샷 예산의
검정은 검출하지 못했다($p=0.11294$). 영과 작은 양수의 구분은 예산에
민감하다.

\section{$P=0$ 품질 감사와 다음 게이트}
기존 기준점 에일리어싱 실험의 $P=0$ 대조군은 정확하게 공통 위상에
불변인데, 초기 32개 Aer 검정에서 명목 5\% 양성이 5개였다.
서로 다른 seed 20개를 추가해 40개 검정 중 양성은 1개였고,
전체는 72개 중 6개였다. 전체 양성 수의 독립 5\% 이항 상측확률은
0.151이다. 추가 실행이 초기 결과를 본 뒤 결정됐으므로 이 전체값은
확증적 $p$값이 아니라 품질 감사의 기술량이다.
36 seed의 raw-count 차이 표준화 평균은 0.0205, 분산은 1.0359였다.
정확 이항확률로 생성한 5000개 대응 패널의 95\% 분산 범위
$[0.9457,1.0540]$, 평균 범위 $[-0.0392,0.0371]$ 안에 들었다.
분산 상측 꼬리비율은 0.0998이다. 초기의 높은 명목 양성 빈도가
추가 seed에서 반복되지 않았으므로 Aer의 체계적 이상이나 물리 신호의
근거로 취급하지 않는다. 원자료와 계산은
\texttt{../results/qiskit\_clock\_alias\_audit.json} 및
\texttt{../results/qiskit\_pzero\_sampling\_audit.json}에 있다.

현재 통과한 것은 \emph{설계된 상태족, $E=\sum k_i^2$,
복귀 검출기, 두 시계 위상, 유한 표지 격자} 안의 개입 식별이다.
더 넓은 ``유한 위상 기준과 시계 아래 관측 차원의 식별 가능성''으로
승격하기 전에는 최소한 다른 검출 POVM과 에너지 스펙트럼,
표지 확률의 불균형, 기준 위상 법칙의 오가정, 위상 기준 자체의 오차에서
같은 판별이 유지되는지 보아야 한다. $P=0$의 정확 null과 Aer 검정력도
계속 분리해 평가한다. QPU 실험은 이 구조 확인 뒤로 보류한다.

\paragraph{재현 자료.}
사전 프로토콜: \texttt{prevalidation\_protocol\_2026-10-04.tex}.
정확 격자: \texttt{../src/relational\_atlas\_i01/holdout/qiskit\_rank\_holdout\_grid.py}와
\texttt{../results/qiskit\_rank\_holdout\_grid.json}.
회로 대조: \texttt{../src/relational\_atlas\_i01/holdout/qiskit\_rank\_holdout\_circuit.py}와
\texttt{../results/qiskit\_rank\_holdout\_circuit.json}.
유한 샷: \texttt{../src/relational\_atlas\_i01/holdout/qiskit\_rank\_holdout\_shots.py}와
\texttt{../results/qiskit\_rank\_holdout\_shots.json}.

\end{document}
